\section{Rotazione dei ruoli}
Per garantire l'acquisizione di competenze trasversali da parte di tutti i membri, la rotazione dei ruoli avverrà secondo le seguenti regole:
\begin{enumerate}
    \item La rotazione è sincronizzata con gli \textbf{sprint bisettimanali}.
    \item Nessun componente può ricoprire lo stesso ruolo per due sprint consecutivi.
    \item Data la presenza di sette componenti e sei ruoli, in ogni sprint sarà presente un ruolo coperto da due persone (ad eccezione del Responsabile, carica ricoperta da un singolo individuo) che avranno compiti distinti. L'identità del ruolo condiviso seguirà le necessità del periodo: fino alla RTB saranno presenti due figure di analista.
    \item Un componente non può effettuare la verifica di prodotti (documenti o codice) da lui stesso realizzati nello sprint corrente o in quello precedente.
    \item I ruoli di analista, progettista e programmatore saranno assegnati secondo una sequenza continuativa: in questo modo chi ha raccolto e interpretato i requisiti di una funzionalità ne conosce il contesto e le motivazioni, e può tradurli in scelte progettuali e implementative coerenti.
    \item Al termine di ogni sprint il Responsabile uscente esegue il passaggio formale di consegne ed aggiorna il consuntivo orario.
    
\end{enumerate}

% ------------------------------------------------------------------------
% 7. CALENDARIO E PERIODI DI RALLENTAMENTO
% ------------------------------------------------------------------------
\section {Calendario e periodi di rallentamento}
Di seguito sono indicate le tappe di riferimento del calendario di progetto:

\vspace{2mm}
\noindent
\begin{tabularx}{\textwidth}{|X c|}
\hline
\textbf{Evento / Milestone} & \textbf{Data prevista} \\
\hline
Aggiudicazione dell'appalto & 11/2026 \\
\hline
Inizio analisi dei requisiti e PoC & 11/2026 \\
\hline
Revisione Requirements and Technology Baseline (RTB) & entro 15/01/2027 \\
\hline
Inizio codifica e progettazione di dettaglio & 18/01/2027 \\
\hline
Revisione Product Baseline (PB) e consegna del prodotto & entro 26/03/2027 \\
\hline
\end{tabularx}

\vspace{2mm}
Il calendario prevede circa 10 settimane per arrivare alla RTB e circa altrettante per la PB. \\
Sono stati individuati due periodi di possibile rallentamento delle attività operative:
\begin{itemize}
    \item \textbf{Vacanze Natalizie (24/12/2026 -- 06/01/2027):} Parziale sospensione delle attività operative per le festività.
    \item \textbf{Sessione d'Esami Invernale (18/01/2027 -- 19/02/2027):} Pianificazione a carico ridotto per consentire lo studio individuale.
\end{itemize}
Per questo motivo si è deciso di collocare la data della presentazione della RTB prima dell'inizio della sessione invernale.
La data di consegna prevista per il progetto è il \textbf{26 marzo 2027}.

% ------------------------------------------------------------------------
% 8. ANALISI DEI RISCHI E MITIGAZIONE
% ------------------------------------------------------------------------
\newpage
\section {Analisi dei rischi e mitigazione}
I principali rischi identificati e le relative strategie di mitigazione sono riportati in tabella:

\vspace{2mm}
\noindent
\begin{tabularx}{\textwidth}{|X c c X|}
\hline
\textbf{Rischio} & \textbf{Probabilità} & \textbf{Impatto} & \textbf{Mitigazione} \\
\hline
Tecnologie nuove per il gruppo & Alta & Medio & Studio approfondito preliminare e realizzazione anticipata di Proof of Concept (PoC). Incontri con esperti di Shopify concordati con l'azienda. \\
\hline
Sottostima delle ore necessarie & Media & Alto & Consuntivo a fine di ogni sprint e rimodulazione tempestiva dei compiti. \\
\hline
Incomprensione dei requisiti & Media & Alto & Incontri periodici di chiarimento con il proponente e verbali di conferma. \\
\hline
Superamento dei limiti dei piani gratuiti & Alta & Alto & Controllo periodico dei consumi, concordare in anticipo con il proponente le alternative disponibili. \\
\hline
Falsi positivi nel motore di rilevamento & Alta & Alto & Sviluppo stratificato, misurare in modo incrementale e mantenere le decisioni in modo deterministico. \\
\hline
Disponibilità ridotta nei periodi d'esame e festività & Alta & Medio & Pianificazione a basso carico nei periodi critici e anticipo delle scadenze chiave. \\
\hline
\end{tabularx}

% ------------------------------------------------------------------------
% 9. CONSEGNA PREVISTA E IMPEGNI FINALI
% ------------------------------------------------------------------------
\section {Consegna prevista e impegni finali}
Il gruppo \textbf{Git Happens} stima di consegnare il prodotto finito entro e non oltre la data del \textbf{26 marzo 2027} (2027/03/26).

Il gruppo si impegna a garantire le 93 ore produttive individuali, a rispettare il preventivo finanziario stabilito di \textbf{\texteuro{} 12.845,00} ed a mantenere un costante coordinamento con il proponente ed il committente.
